% =============================================================================
% FOM Seminararbeit – Wirtschaftsinformatik / IT-Recht
% Template basierend auf FOM-Leitfaden V1.4 (März 2024)
% und Jäger/Kümpel/Seng Leitfaden (Januar 2024)
% =============================================================================
\documentclass[12pt,oneside,titlepage,listof=totoc,bibliography=totoc]{scrartcl}

% ---------------------------------------------------------------------------
% Formalia -- MUSS als erstes stehen
% ---------------------------------------------------------------------------
% Setzt Raender, Schrift, Zitierstil und die leitfadenspezifischen Sonder-
% regeln. Einstellen tust du das in formalia/konfig.tex, nicht hier.
% Diese Zeile muss vor fontspec, geometry und biblatex stehen -- deren
% Optionen lassen sich nachtraeglich nicht mehr aendern.
% =============================================================================
% FORMALIA-MASCHINERIE -- hier musst du nichts aendern
% =============================================================================
% Diese Datei wird als erstes nach \documentclass geladen. Sie muss VOR
% fontspec, geometry und biblatex laufen, weil Schrift, Raender und Zitierstil
% Ladeoptionen dieser Pakete sind und sich nachtraeglich nicht mehr aendern
% lassen. Genau daran ist ein frueherer Entwurf gescheitert, der die Schalter
% in skripte/meta.tex legte -- die wird erst kurz vor \begin{document} gelesen.
%
% Reihenfolge: Standardwerte -> Leitfadenprofil -> Overlay -> eigene Vorgaben.
% Was spaeter kommt, gewinnt.
% =============================================================================

% =============================================================================
% FORMALIA-SCHALTER -- hier stellst du deinen Leitfaden ein
% =============================================================================
% An der FOM gibt es nicht DEN einen Leitfaden. Welcher fuer dich gilt, steht
% in deiner Modulbeschreibung bzw. sagen dir deine Pruefenden. Stelle hier ein,
% welcher gilt -- die Vorlage setzt Raender, Schrift, Zitierstil und die
% leitfadenspezifischen Sonderregeln dann automatisch richtig.
%
% Hintergrund zu jeder Einstellung: ../anleitung/02_formalia.md
%
% WICHTIG: Individuelle Vorgaben deiner Pruefenden haben IMMER Vorrang vor dem
% Leitfaden. Solche Abweichungen traegst du nicht hier ein, sondern in
% formalia/profil_eigen.tex -- dann bleibt nachvollziehbar, was Leitfaden ist
% und was Pruefervorgabe.
% =============================================================================

% --- Welcher Leitfaden gilt? -------------------------------------------------
%   wi   = Gestaltung wissenschaftlicher Arbeiten, V1.4 (03/2024)
%          Wirtschaftsinformatik & Ingenieurwesen
%   jks  = Formale Gestaltung, Jaeger/Kuempel/Seng (01/2024)
%          allgemeiner FOM-Standard (Wirtschaft & Recht u. a.)
\def\FormaliaLeitfaden{jks}

% --- Zitierstil --------------------------------------------------------------
%   fussnote    = Kurzbeleg in der Fussnote  (JKS-Standard, WI zulaessig)
%   authoryear  = Harvard/Autor-Jahr im Text (JKS zulaessig, WI zulaessig)
%   ieee        = numerisch [12]             (nur WI, in der Informatik ueblich)
%
% ACHTUNG: Diesen Wert einmal am Anfang festlegen und dann NICHT mehr wechseln.
% Der Schalter aendert nur die Ausgabe, nicht den Sinn der bereits im Text
% gesetzten \autocite-Befehle. Ein spaeter Wechsel von fussnote auf authoryear
% verschiebt saemtliche Belege aus den Fussnoten in den Fliesstext -- das musst
% du dann sprachlich nacharbeiten.
\def\FormaliaZitierstil{fussnote}

% --- Schrift -----------------------------------------------------------------
%   tnr    = Times New Roman 12 pt (Standard beider Leitfaeden)
%            Gesetzt wird TeX Gyre Termes, ein metrisch identischer freier Klon.
%   arial  = Arial als zulaessige Alternative (WI 11 pt / JKS 11,5 pt)
%            Gesetzt wird TeX Gyre Heros, ein metrisch identischer freier Klon.
\def\FormaliaSchrift{tnr}

% --- Optionales Overlay ------------------------------------------------------
%   ifes   = zusaetzliche Vorgaben fuer empirische Arbeiten (Gansser u. a.)
%   kein   = kein Overlay
% Das Overlay ersetzt den Leitfaden nicht, es ergaenzt ihn.
\def\FormaliaOverlay{kein}


% --- Standardwerte (gelten, bis ein Profil sie ueberschreibt) ----------------
\def\formaliaLeitfadenName{--}
\def\formaliaLeitfadenStand{--}
\def\formaliaGeprueftAm{--}

\def\formaliaRandLinks{4cm}
\def\formaliaRandRechts{2cm}
\def\formaliaRandOben{4cm}
\def\formaliaRandUnten{2cm}

\def\formaliaSchriftHaupt{TeX Gyre Termes}
\def\formaliaSchriftGroesse{12pt}

\def\formaliaSeitenzahlPos{C}          % C = oben mittig, R = oben rechts
\def\formaliaBibOpt{style=ext-authoryear-ibid,autocite=footnote,%
  mergedate=false,innamebeforetitle,dashed=false}

\newif\ifformaliaQuelleEigen           % "Quelle: Eigene Darstellung" verlangt?
\newif\ifformaliaSperrvermerkImTOC     % Sperrvermerk im Inhaltsverzeichnis?
\newif\ifformaliaInternetquellenSeparat

% --- Leitfadenprofil laden ---------------------------------------------------
% Bewusst als explizite Fallunterscheidung statt \input{...\FormaliaLeitfaden}:
% lesbarer, und ein Tippfehler im Schalter liefert eine klare Fehlermeldung
% statt "File not found".
\def\formaliaInternwi{wi}
\def\formaliaInternjks{jks}
\ifx\FormaliaLeitfaden\formaliaInternwi
  % =============================================================================
% PROFIL: "Gestaltung wissenschaftlicher Arbeiten", V1.4, Maerz 2024.
%         Wirtschaftsinformatik & Ingenieurwesen.
% -----------------------------------------------------------------------------
% Leitfaden-Fassung : V1.4, Maerz 2024
% Ausgewertet am    : 2026-08-26
% Belege            : ../anleitung/02_formalia.md
%
% Leitfaeden werden ueberarbeitet. Hol dir die aktuelle Fassung aus dem
% FOM Online-Campus und pruefe die Werte hier dagegen, bevor du abgibst.
% =============================================================================

\def\formaliaLeitfadenName{dem FOM-Leitfaden Wirtschaftsinformatik}
\def\formaliaLeitfadenStand{V1.4, Maerz 2024}
\def\formaliaGeprueftAm{2026-08-26}

% WI nennt die Masse "ungefaehr". Die Vorlage trifft die strengere Lesart.
\def\formaliaRandLinks{4cm}
\def\formaliaRandRechts{2cm}
\def\formaliaRandOben{4cm}
\def\formaliaRandUnten{2cm}

\def\formaliaSeitenzahlPos{C}

% WI verbietet den Vermerk "Quelle: Eigene Darstellung" ausdruecklich.
\formaliaQuelleEigenfalse

% WI: Sperrvermerk ohne Seitenzahl und NICHT im Inhaltsverzeichnis.
\formaliaSperrvermerkImTOCfalse

% WI trennt Internetquellen nicht ab.
\formaliaInternetquellenSeparatfalse

\else
  \ifx\FormaliaLeitfaden\formaliaInternjks
    % =============================================================================
% PROFIL: Jaeger/Kuempel/Seng -- "Formale Gestaltung wissenschaftlicher
%         Arbeiten", Januar 2024. Allgemeiner FOM-Standard.
% -----------------------------------------------------------------------------
% Leitfaden-Fassung : Januar 2024
% Ausgewertet am    : 2026-08-26
% Belege            : ../anleitung/02_formalia.md
%
% Leitfaeden werden ueberarbeitet. Hol dir die aktuelle Fassung aus dem
% FOM Online-Campus und pruefe die Werte hier dagegen, bevor du abgibst.
% =============================================================================

\def\formaliaLeitfadenName{dem FOM-Leitfaden von Jaeger, Kuempel und Seng}
\def\formaliaLeitfadenStand{Januar 2024}
\def\formaliaGeprueftAm{2026-08-26}

% Seitenlayout: der breite linke Rand ist Korrektur- und Bindungsrand.
\def\formaliaRandLinks{4cm}
\def\formaliaRandRechts{2cm}
\def\formaliaRandOben{4cm}
\def\formaliaRandUnten{2cm}

% Seitenzahl oben. JKS laesst mittig und "oben rechts am Textrand" zu.
\def\formaliaSeitenzahlPos{C}

% JKS verlangt bei selbst erstellten Abbildungen/Tabellen den Vermerk
% "Quelle: Eigene Darstellung". (WI verbietet ihn -- echter Widerspruch.)
\formaliaQuelleEigentrue

% JKS: Sperrvermerk mit roemischer Seitenzahl UND im Inhaltsverzeichnis.
\formaliaSperrvermerkImTOCtrue

% JKS: Internetquellen kommen separat ans Ende des Literaturverzeichnisses.
\formaliaInternetquellenSeparattrue

  \else
    \errmessage{[formalia] Unbekannter Leitfaden in formalia/konfig.tex.
      Erlaubt sind: wi, jks}
  \fi
\fi

% --- Zitierstil --------------------------------------------------------------
\def\formaliaInternfussnote{fussnote}
\def\formaliaInternauthoryear{authoryear}
\def\formaliaInternieee{ieee}
\ifx\FormaliaZitierstil\formaliaInternfussnote
  \def\formaliaBibOpt{style=ext-authoryear-ibid,autocite=footnote,%
    mergedate=false,innamebeforetitle,dashed=false}
\else
  \ifx\FormaliaZitierstil\formaliaInternauthoryear
    \def\formaliaBibOpt{style=authoryear,autocite=inline,%
      mergedate=false,dashed=false}
  \else
    \ifx\FormaliaZitierstil\formaliaInternieee
      \ifx\FormaliaLeitfaden\formaliaInternjks
        \errmessage{[formalia] Zitierstil ieee ist nach JKS nicht zulaessig.
          JKS erlaubt Chicago in Fussnoten oder Harvard im Text.}
      \fi
      \def\formaliaBibOpt{style=numeric-comp,autocite=inline,sorting=none}
    \else
      \errmessage{[formalia] Unbekannter Zitierstil in formalia/konfig.tex.
        Erlaubt sind: fussnote, authoryear, ieee}
    \fi
  \fi
\fi

% --- Schrift -----------------------------------------------------------------
\def\formaliaInterntnr{tnr}
\def\formaliaInternarial{arial}
\ifx\FormaliaSchrift\formaliaInterntnr
  \def\formaliaSchriftHaupt{TeX Gyre Termes}
  \def\formaliaSchriftGroesse{12pt}
\else
  \ifx\FormaliaSchrift\formaliaInternarial
    \def\formaliaSchriftHaupt{TeX Gyre Heros}
    % WI 11 pt, JKS 11,5 pt -- KOMA kennt keine halben Punkt als Klassenoption,
    % deshalb wird 11.5pt ueber \KOMAoptions gesetzt.
    \ifx\FormaliaLeitfaden\formaliaInternwi
      \def\formaliaSchriftGroesse{11pt}
    \else
      \def\formaliaSchriftGroesse{11.5pt}
    \fi
  \else
    \errmessage{[formalia] Unbekannte Schrift in formalia/konfig.tex.
      Erlaubt sind: tnr, arial}
  \fi
\fi

% --- Overlay -----------------------------------------------------------------
\def\formaliaInternifes{ifes}
\ifx\FormaliaOverlay\formaliaInternifes
  % =============================================================================
% OVERLAY: ifes -- Leitfaden fuer empirische Arbeiten (Gansser u. a.)
% -----------------------------------------------------------------------------
% Ein Overlay ersetzt den Leitfaden NICHT, es ergaenzt ihn. Geladen wird es
% zusaetzlich zu profil_wi bzw. profil_jks.
%
% Ausgewertet am : 2026-08-26
% Beleg          : ../anleitung/02_formalia.md
% =============================================================================

% ifes laesst den rechten Rand bis auf 1 cm zu. Die Vorlage bleibt bewusst bei
% 2 cm -- schmaler ist zulaessig, aber nicht besser lesbar. Wer den Platz
% braucht, setzt in profil_eigen.tex \def\formaliaRandRechts{1cm}.

% Inhaltliche ifes-Vorgabe, die LaTeX nicht pruefen kann und die deshalb in
% deiner FORMALIA.md stehen muss:
%   Der empirische Teil muss mehr als 50 % der Arbeit ausmachen.

\fi

% --- Eigene Vorgaben (Pruefende schlagen Leitfaden) --------------------------
\InputIfFileExists{formalia/profil_eigen}{}{}

% --- Werte an die Pakete uebergeben ------------------------------------------
% \edef expandiert die Makros JETZT, damit die Pakete fertige Werte sehen.
\edef\formaliaInterngeo{a4paper,%
  left=\formaliaRandLinks,right=\formaliaRandRechts,%
  top=\formaliaRandOben,bottom=\formaliaRandUnten,%
  headheight=14pt,headsep=1cm}
\expandafter\PassOptionsToPackage\expandafter{\formaliaInterngeo}{geometry}
\expandafter\PassOptionsToPackage\expandafter{\formaliaBibOpt}{biblatex}

\edef\formaliaInternfs{\formaliaSchriftGroesse}
\KOMAoptions{fontsize=\formaliaInternfs}

% --- Helfer fuer den Fliesstext ----------------------------------------------
% Unter selbst erstellten Abbildungen/Tabellen verwenden. JKS verlangt den
% Vermerk, WI verbietet ihn ausdruecklich -- dieser Befehl macht das Richtige.
\newcommand{\formaliaQuellenvermerk}{%
  \ifformaliaQuelleEigen Quelle: Eigene Darstellung\fi}

% Fuer die Methodik-Passage der Einleitung: den gewaehlten Leitfaden nennen
% ist selbst eine Formvorgabe.
\newcommand{\formaliaLeitfadenSatz}{%
  Die formale Gestaltung folgt \formaliaLeitfadenName\ (\formaliaLeitfadenStand).}


% ---------------------------------------------------------------------------
% Sprache
% ---------------------------------------------------------------------------
\usepackage[ngerman]{babel}
\usepackage[babel,german=quotes]{csquotes}

% ---------------------------------------------------------------------------
% Schriftart -- der Name kommt aus dem Formalia-Profil (TeX Gyre Termes als
% metrisch identischer Times-New-Roman-Klon bzw. TeX Gyre Heros fuer Arial).
% ---------------------------------------------------------------------------
\usepackage{fontspec}
\expandafter\setmainfont\expandafter{\formaliaSchriftHaupt}

% ---------------------------------------------------------------------------
% Seitengeometrie -- die Masse kommen aus dem Formalia-Profil
% (formalia/profil_*.tex). Der breite linke Rand ist Korrektur- und
% Bindungsrand, deshalb ist er kein Gestaltungsspielraum.
\usepackage{geometry}

% ---------------------------------------------------------------------------
% Zeilenabstand 1,5-zeilig (Leitfaden S. 16)
% ---------------------------------------------------------------------------
\usepackage{setspace}
\onehalfspacing

% ---------------------------------------------------------------------------
% Absatzformatierung (Leitfaden: vor 0pt, nach 6pt, Blocksatz)
% ---------------------------------------------------------------------------
\setlength{\parindent}{0mm}
\setlength{\parskip}{6pt plus 2pt minus 1pt}

% ---------------------------------------------------------------------------
% Kopf-/Fusszeile -- Position der Seitenzahl kommt aus dem Formalia-Profil
% ---------------------------------------------------------------------------
\usepackage{fancyhdr}
\pagestyle{fancy}
\fancyhf{}
\edef\formaliaInternkopf{\noexpand\fancyhead[\formaliaSeitenzahlPos]{\noexpand\thepage}}
\formaliaInternkopf
\renewcommand{\headrulewidth}{0pt}

% ---------------------------------------------------------------------------
% Grafik & Farben
% ---------------------------------------------------------------------------
\usepackage{graphicx}
\graphicspath{{./abbildungen/}}
\usepackage[table]{xcolor}
% floatrow laedt float selbst und stellt damit die Option [H] bereit
% (exakte Platzierung). Ein zusaetzliches \usepackage{float} bricht den
% Build ab: "Do not use float package with floatrow."
\usepackage[capposition=top]{floatrow}

% ---------------------------------------------------------------------------
% Beschriftungen (Leitfaden: linksbündig, fett)
% ---------------------------------------------------------------------------
\usepackage[
  singlelinecheck=false,
  labelfont=bf,
  font=bf,
  aboveskip=12pt,
  belowskip=6pt
]{caption}

% ---------------------------------------------------------------------------
% Tabellen
% ---------------------------------------------------------------------------
\usepackage{array}
\usepackage{tabularx}
\usepackage{multirow}
\usepackage{booktabs}

% ---------------------------------------------------------------------------
% Listen
% ---------------------------------------------------------------------------
\usepackage{enumitem}

% ---------------------------------------------------------------------------
% Mathematik & Symbole
% ---------------------------------------------------------------------------
\usepackage{amsmath}
\usepackage{amssymb}

% ---------------------------------------------------------------------------
% Quellcode-Listings
% ---------------------------------------------------------------------------
\usepackage{listings}
\lstset{
  numbers=left,
  numberstyle=\tiny,
  numbersep=5pt,
  breaklines=true,
  showstringspaces=false,
  frame=l,
  xleftmargin=5pt,
  xrightmargin=5pt,
  basicstyle=\ttfamily\scriptsize,
  stepnumber=1,
  commentstyle=\color{gray},
  keywordstyle=\color{blue},
  stringstyle=\color{purple}
}

% ---------------------------------------------------------------------------
% Fußnoten (einzeilig, 10pt, Leitfaden S. 18)
% ---------------------------------------------------------------------------
\usepackage[hang,multiple]{footmisc}
\setlength{\footnotemargin}{1em}

% ---------------------------------------------------------------------------
% Silbentrennung (Leitfaden S. 17)
% ---------------------------------------------------------------------------
\usepackage{microtype}
\hyphenpenalty=5000
\exhyphenpenalty=5000
\hbadness=10000
\vbadness=10000
\hfuzz=40pt
\vfuzz=40pt

% ---------------------------------------------------------------------------
% Abkürzungsverzeichnis
% ---------------------------------------------------------------------------
\usepackage[printonlyused]{acronym}

% ---------------------------------------------------------------------------
% Verhinderung von Schusterjungen/Hurenkinder
% ---------------------------------------------------------------------------
\usepackage[all]{nowidow}
\setlength{\marginparwidth}{2cm}

% ---------------------------------------------------------------------------
% Literaturverzeichnis – BibLaTeX/Biber
% Stil: ext-authoryear-ibid (FOM 2018/2024 Fußnoten-Stil)
% ---------------------------------------------------------------------------
\usepackage{url}
\urlstyle{same}
\usepackage{xurl}
\usepackage[
  backend=biber,
  % style und autocite kommen aus dem Formalia-Profil (siehe formalia/laden.tex).
  % Beides sind Ladeoptionen und liessen sich hier nicht mehr umstellen.
  maxcitenames=3,
  maxbibnames=999,
  date=iso,
  seconds=true,
  urldate=iso,
  doi=false,
  useprefix=true,
  mincrossrefs=1
]{biblatex}
\addbibresource{literatur/literatur.bib}
\setcounter{biburllcpenalty}{7000}
\setcounter{biburlucpenalty}{8000}
\setcounter{biburlnumpenalty}{9000}

% Einzeiliger Zeilenabstand im Literaturverzeichnis (Leitfaden S. 14)
\AtBeginBibliography{\singlespacing\sloppy}

% Anpassungen für dt. BibLaTeX (et al. statt u. a.)
\DefineBibliographyStrings{ngerman}{
  andothers = {{et\,al\adddot}},
}

% Lade FOM-spezifische BibLaTeX-Modifikationen
% =============================================================================
% BibLaTeX-Modifikationen für FOM 2018/2024 Zitationsstil
% Anpassungen für ext-authoryear-ibid
% =============================================================================

% Nachnamen in Kapitälchen im Literaturverzeichnis (optional)
% \renewcommand*{\mkbibnamefamily}[1]{\textsc{#1}}

% Doppelpunkt nach Autor statt Punkt
\renewcommand*{\labelnamepunct}{\addcolon\space}

% "Vgl." als Standard-Prenote
\renewcommand{\mkccitation}[1]{ #1}

% URL-Datum-Präfix
\DeclareFieldFormat{urldate}{(Zugriff am: #1)}

% Online-Quellen: URL anzeigen
\DeclareFieldFormat[online]{url}{\url{#1}}

% Gerichtsentscheidungen: Datum und Aktenzeichen stehen im Titel ("Urteil vom
% 09.09.2021 -- I ZR 113/20"), das Jahr im Label. Der Stil haengt sonst das
% ISO-Datum noch einmal ans Ende des Eintrags. Das Label bleibt erhalten,
% weil Biber es vor diesem Hook berechnet.
\AtEveryBibitem{%
  \ifentrytype{jurisdiction}{%
    \clearfield{date}\clearfield{year}\clearfield{month}\clearfield{day}%
  }{}%
}


% ---------------------------------------------------------------------------
% Hyperlinks (PDF-Metadaten)
% ---------------------------------------------------------------------------
\usepackage[hyperfootnotes=false]{hyperref}
\hypersetup{
  colorlinks=true,
  breaklinks=true,
  linkcolor=black,
  citecolor=black,
  menucolor=black,
  urlcolor=black,
  linktoc=all,
  bookmarksnumbered=false,
  pdfpagemode=UseOutlines,
}

% ---------------------------------------------------------------------------
% Metadaten (werden in skripte/meta.tex definiert)
% ---------------------------------------------------------------------------
% =============================================================================
% Meta-Informationen der Arbeit
% -----------------------------------------------------------------------------
% HIER trägst du als Erstes deine persönlichen Daten ein. Diese Werte tauchen
% automatisch auf dem Deckblatt und in den PDF-Eigenschaften auf. Du musst sonst
% nirgends im Projekt etwas an deinen Daten ändern.
%
% Tipp: Ersetze nur den Text INNERHALB der geschweiften Klammern { ... }.
% =============================================================================

% --- Titel der Arbeit ---
\newcommand{\myTitel}{Mustertitel der wissenschaftlichen Arbeit}
\newcommand{\myUntertitel}{Ein aussagekräftiger Untertitel (optional)}

% --- Autor ---
\newcommand{\myAutor}{Vorname Nachname}
\newcommand{\myMatrikelnummer}{000000}

% --- Studiengang ---
\newcommand{\myStudiengang}{Wirtschaftsinformatik}
\newcommand{\myStudienort}{Musterstadt}

% --- Betreuung ---
\newcommand{\myBetreuer}{Prof.~Dr.~Mustername}
\newcommand{\myZweitbetreuer}{} % leer lassen, wenn nicht vorhanden

% --- Datum ---
\newcommand{\myAbgabedatum}{01.01.2026}

% --- Art der Arbeit ---
% z.B. Hausarbeit, Seminararbeit, Projektarbeit, Bachelorarbeit, Masterarbeit
\newcommand{\myArbeit}{Seminararbeit}

% --- Wortzahl (optional) ---
% Manche Pruefende verlangen die Wortzahl auf dem Titelblatt. Leer lassen,
% wenn nicht gefordert -- dann erscheint die Zeile nicht. Zaehlen so, wie
% deine Pruefenden zaehlen: ./skripte/woerter.sh <erste> <letzte Seite>
\newcommand{\myWortzahl}{}

% --- Modul / Kurs ---
\newcommand{\myModul}{Modulname}

% --- Hochschule ---
\newcommand{\myHochschule}{FOM Hochschule für Oekonomie \& Management}

% --- Firma (nur bei Arbeiten mit Sperrvermerk) ---
% Leer lassen, wenn deine Arbeit keinen Sperrvermerk braucht.
\newcommand{\myFirma}{Musterfirma GmbH}

% =============================================================================
% Textbausteine fuer haeufige Abkuerzungen
% =============================================================================
% Warum das wichtig ist: "z. B." wird mit einem SCHMALEN, GESCHUETZTEN
% Leerzeichen gesetzt. Schreibst du "z. B." mit normalem Leerzeichen, darf
% LaTeX zwischen "z." und "B." umbrechen -- ein klassischer Typografiefehler,
% der in der Korrektur auffaellt. Schreibst du "z.B." ohne Leerzeichen, ist es
% ebenfalls falsch.
%
% Benutzung im Fliesstext:   \zb  \dah  \ua  \vgl  \bzw  \ca
% Vor einem Komma oder Punkt die OL-Variante (ohne folgendes Leerzeichen):
%   ... \zbol, ...
%
% Uebernommen aus dem FOM-LaTeX-Template von Andy Grunwald (MIT) und auf die
% an der FOM gebraeuchlichen Formen reduziert.
% =============================================================================

\newcommand{\zbol}{z.\,B.}
\newcommand{\zb}{\zbol\ }
\newcommand{\dahol}{d.\,h.}  % \dh ist von LaTeX belegt (Eth-Zeichen), daher \dah
\newcommand{\dah}{\dahol\ }
\newcommand{\uaol}{u.\,a.}
\newcommand{\ua}{\uaol\ }
\newcommand{\vglol}{Vgl.}
\newcommand{\vgl}{\vglol\ }
\newcommand{\bzwol}{bzw.}
\newcommand{\bzw}{\bzwol\ }
\newcommand{\caol}{ca.}
\newcommand{\ca}{\caol\ }
\newcommand{\ggfol}{ggf.}
\newcommand{\ggf}{\ggfol\ }
\newcommand{\os}{o.\,S.}      % ohne Seitenangabe (Online-Quellen)
\newcommand{\oj}{o.\,J.}      % ohne Jahr
\newcommand{\ov}{o.\,V.}      % ohne Verfasser


% ---------------------------------------------------------------------------
% Glossar (optional)
% ---------------------------------------------------------------------------
\usepackage{glossaries}
\glstoctrue
\makenoidxglossaries
% =============================================================================
% Glossar-Einträge
% =============================================================================

% Beispiel:
% \newglossaryentry{informationssystem}{
%   name={Informationssystem},
%   description={Ein soziotechnisches System bestehend aus Menschen,
%   Informations- und Kommunikationstechnik und Organisation}
% }


% ---------------------------------------------------------------------------
% Anhang-Paket
% ---------------------------------------------------------------------------
\usepackage{appendix}

% ---------------------------------------------------------------------------
% Weitere Gliederungsebene (paragraph als nummerierte Ebene)
% ---------------------------------------------------------------------------
\setcounter{secnumdepth}{4}
\setcounter{tocdepth}{4}

% ============================= DOKUMENT START ==============================
\begin{document}

\pagenumbering{Roman}

% --- Titelseite ---
% =============================================================================
% Titelseite (FOM-Leitfaden konform)
% =============================================================================
\begin{titlepage}
  \begin{center}

    \vspace*{1cm}

    {\large \myHochschule}\\[0.5cm]
    {\large Studiengang: \myStudiengang}\\[0.5cm]
    {\large Studienort: \myStudienort}\\[2cm]

    {\LARGE \textbf{\myTitel}}\\[0.5cm]
    {\large \myUntertitel}\\[2cm]

    {\large \myArbeit}\\[0.5cm]
    {\large im Modul \myModul}\\[2cm]

    \begin{tabular}{rl}
      \textbf{Verfasser/in:}    & \myAutor \\
      \textbf{Matrikelnummer:}   & \myMatrikelnummer \\
      \textbf{Erstbetreuer/in:}  & \myBetreuer \\
      \ifx\myWortzahl\empty\else
      \textbf{Wortzahl:}         & \myWortzahl \\
      \fi
    \end{tabular}

    \vfill

    {\large \myAbgabedatum}

  \end{center}
\end{titlepage}


% --- Inhaltsverzeichnis ---
\setcounter{page}{2}
\tableofcontents
\newpage

% --- Abbildungsverzeichnis ---
% Erst entkommentieren (% entfernen), wenn die Arbeit Abbildungen enthaelt.
%\listoffigures
%\newpage

% --- Tabellenverzeichnis ---
% Erst entkommentieren, wenn die Arbeit Tabellen enthaelt.
%\listoftables
%\newpage

% --- Abkürzungsverzeichnis ---
\addcontentsline{toc}{section}{Abkürzungsverzeichnis}
% =============================================================================
% Abkürzungsverzeichnis
% -----------------------------------------------------------------------------
% Nur FACHSPEZIFISCHE Abkürzungen aufnehmen. Triviale wie "z.B." oder "usw."
% gehören NICHT hierher. Alphabetisch sortiert halten.
%
% Im Text nutzt du \ac{KUERZEL}: Bei der ersten Verwendung wird der Begriff
% automatisch ausgeschrieben (inkl. Kürzel), danach erscheint nur das Kürzel.
% Dank [printonlyused] erscheinen hier NUR tatsächlich verwendete Abkürzungen.
% =============================================================================
\section*{Abkürzungsverzeichnis}

\begin{acronym}[DSGVO]  % längstes Kürzel als Referenzbreite für die Einrückung
  \acro{BSI}{Bundesamt für Sicherheit in der Informationstechnik}
  \acro{DSGVO}{Datenschutz-Grundverordnung}
  \acro{KI}{Künstliche Intelligenz}
  \acro{KMU}{kleine und mittlere Unternehmen}
  \acro{LLM}{Large Language Model}
\end{acronym}

\newpage

% --- Glossar (optional) ---
% \printnoidxglossaries
% \newpage

% ============================= TEXTTEIL ====================================
\pagenumbering{arabic}
\setcounter{page}{1}

% Kapitel einbinden
% -------------------------------------------------------------------------
% ACHTUNG: Die folgenden vier Kapitel sind nur ein BEISPIELGERUEST, kein
% Pflicht-Schema. Passe sie an deine Forschungsfrage an -- loesche, benenne
% um oder ergaenze Kapitel frei. (Strukturhilfe: anleitung/06_gliederung.md)
%
% Jede Datei unter kapitel/ ist EIN Kapitel. Neues Kapitel anlegen:
%   1. Datei kapitel/05_xyz.tex erstellen (mit \section{...} beginnen)
%   2. Hier eine Zeile \input{kapitel/05_xyz} ergaenzen
% Die Reihenfolge HIER bestimmt die Reihenfolge im Dokument.
% -------------------------------------------------------------------------
% =============================================================================
% Kapitel 1: Einleitung
% -----------------------------------------------------------------------------
% Die Einleitung fuehrt in das Thema ein und endet mit Forschungsfrage,
% Methodik und Aufbau. Sie ist das einzige Kapitel OHNE Kapitel-Einleitungssatz
% vor der ersten \subsection (alle anderen Kapitel brauchen 2--4 einleitende
% Saetze zwischen \section und erster \subsection -- siehe FOM-Leitfaden).
% =============================================================================
\section{Einleitung}
\label{sec:einleitung}

% --- 1.1 Problemstellung und Relevanz ---
\subsection{Problemstellung und Relevanz}
\label{subsec:problemstellung}

Hier beschreibst du das Problem und warum es relevant ist. Beginne mit einem
konkreten Aufhänger und arbeite dich zur wissenschaftlichen Lücke vor. Jede
Tatsachenbehauptung wird belegt. Ein indirektes (sinngemäßes) Zitat sieht so
aus:\autocite[Vgl.][S.~15]{mustermann2024}
Ein direktes Zitat steht in Anführungszeichen und ohne \enquote{Vgl.}:
\enquote{Wörtlich übernommener Satz.}\autocite[][S.~16]{mustermann2024}

Fachbegriffe werden bei der ersten Verwendung eingeführt. Danach genügt die
Abkürzung, z.\,B. \ac{KMU}. Bei der ersten Nennung schreibt \verb|\ac| den
Begriff automatisch aus.

% --- 1.2 Forschungsfrage ---
\subsection{Forschungsfrage}
\label{subsec:forschungsfrage}

Aus der Problemstellung ergibt sich die folgende Forschungsfrage:

\begin{quote}
  \itshape
  Hier steht deine eine, präzise Forschungsfrage. Sie ist spezifisch,
  beantwortbar und im Umfang der Arbeit zu bewältigen.
\end{quote}

% --- 1.3 Methodik und Aufbau ---
\subsection{Methodik und Aufbau der Arbeit}
\label{subsec:methodik}

Hier erläuterst du, mit welcher Methode du die Frage beantwortest
(z.\,B. Literaturanalyse, empirische Erhebung, Fallstudie) und wie die Arbeit
aufgebaut ist. Verweise auf die folgenden Kapitel mit \verb|\ref|:
Kapitel~\ref{sec:grundlagen} legt die Grundlagen, Kapitel~\ref{sec:hauptteil}
analysiert den Kern der Fragestellung, und Kapitel~\ref{sec:fazit} fasst die
Ergebnisse zusammen.

Am Ende dieses Abschnitts nennst du den verwendeten FOM-Leitfaden, an dem sich
die formale Gestaltung orientiert (FOM-Regel).

\newpage
% =============================================================================
% Kapitel 2: Grundlagen
% -----------------------------------------------------------------------------
% Hier werden die Begriffe und Konzepte geklärt, die der Hauptteil voraussetzt.
% Beachte: Zwischen \section und der ersten \subsection stehen 2--4 Sätze
% (Kapitel-Einleitung), die den Inhalt zusammenfassen und an Kapitel 1 anknüpfen.
% =============================================================================
\section{Grundlagen}
\label{sec:grundlagen}

Dieses Kapitel klärt die zentralen Begriffe der Arbeit und ordnet sie in den
Forschungsstand ein. Es schafft damit die Grundlage für die Analyse in
Kapitel~\ref{sec:hauptteil}. Behandelt werden zunächst die begrifflichen
Grundlagen und anschließend der relevante Stand der Forschung.

% --- 2.1 ---
\subsection{Begriffliche Grundlagen}
\label{subsec:begriffe}

Definiere hier die Schlüsselbegriffe. Eine Definition wird belegt:
Ein \ac{KMU} ist nach gängiger Definition \dots\autocite[Vgl.][S.~3]{musterinstitut2023}

% --- 2.2 ---
\subsection{Stand der Forschung}
\label{subsec:forschungsstand}

Hier ordnest du die wichtigsten Quellen ein und zeigst, wo deine Arbeit
ansetzt. So bindest du eine Abbildung ein (Datei liegt unter abbildungen/):

% Beispiel fuer eine Abbildung -- erst einkommentieren, wenn ein Bild vorliegt,
% und im main.tex \listoffigures aktivieren.
% \begin{figure}[H]
%   \centering
%   \includegraphics[width=0.8\textwidth]{beispielbild.png}
%   \caption{Beschreibung der Abbildung}
%   \label{fig:beispiel}
% \end{figure}
% Im Text immer referenzieren: Abbildung~\ref{fig:beispiel} zeigt \dots

So bindest du eine Tabelle ein:

\begin{table}[H]
  \centering
  \caption{Beispielhafte Gegenüberstellung}
  \label{tab:beispiel}
  \begin{tabularx}{\textwidth}{@{}lX@{}}
    \toprule
    \textbf{Kriterium} & \textbf{Beschreibung} \\
    \midrule
    Merkmal A & Erläuterung zu Merkmal A \\
    Merkmal B & Erläuterung zu Merkmal B \\
    \bottomrule
  \end{tabularx}
\end{table}

\noindent Tabelle~\ref{tab:beispiel} stellt die Kriterien gegenüber.

\newpage
% =============================================================================
% Kapitel 3: Hauptteil / Analyse
% -----------------------------------------------------------------------------
% Der Kern der Arbeit. Hier wird die Forschungsfrage bearbeitet. Pro
% Unterkapitel eine zentrale Aussage, im Dreischritt: Begriff/Gegenstand klären
% -- entfalten -- Folgerung ziehen und überleiten.
% =============================================================================
\section{Analyse}
\label{sec:hauptteil}

Aufbauend auf den Grundlagen aus Kapitel~\ref{sec:grundlagen} bearbeitet dieses
Kapitel die Forschungsfrage. Es entwickelt das Argument schrittweise und führt
die Teilergebnisse zusammen.

% --- 3.1 ---
\subsection{Erster Analyseschritt}
\label{subsec:analyse1}

Jeder Absatz trägt genau einen Gedanken und ist belegt. Mehrere Quellen lassen
sich kombinieren.\autocite[Vgl.][S.~22]{mustermann2024}\autocite[Vgl.][S.~5]{musterinstitut2023}
Eine Online-Quelle ohne Seitenzahl wird mit Absatz oder Randnummer
zitiert.\autocite[Vgl.][o.\,S.]{musterbehoerde2025}

% --- 3.2 ---
\subsection{Zweiter Analyseschritt}
\label{subsec:analyse2}

Der letzte Absatz eines Unterkapitels zieht eine Folgerung und leitet zum
nächsten über, ohne dessen Ergebnis vorwegzunehmen.

\newpage
% =============================================================================
% Kapitel 4: Fazit und Ausblick
% -----------------------------------------------------------------------------
% Das Fazit beantwortet die Forschungsfrage aus der Einleitung explizit.
% Keine neuen Quellen, keine neuen Argumente -- nur Zusammenführung. Der
% Ausblick benennt offene Fragen.
% =============================================================================
\section{Fazit und Ausblick}
\label{sec:fazit}

% --- 4.1 ---
\subsection{Beantwortung der Forschungsfrage}
\label{subsec:beantwortung}

Hier beantwortest du die in Kapitel~\ref{subsec:forschungsfrage} gestellte
Forschungsfrage direkt und fasst die tragenden Ergebnisse aus
Kapitel~\ref{sec:hauptteil} zusammen. Prüfe: Beantwortet dieser Abschnitt
wirklich die Frage, die du zu Beginn gestellt hast?

% --- 4.2 ---
\subsection{Ausblick}
\label{subsec:ausblick}

Welche Fragen bleiben offen? Wo könnte weitere Forschung ansetzen? Der Ausblick
ist kurz und konkret, keine Wiederholung des Fazits.


% ============================= NACHSPANN ===================================
\newpage

% --- Literaturverzeichnis ---
\printbibliography[nottype=online,heading=bibintoc,title={Literaturverzeichnis}]
\newpage
\printbibliography[type=online,heading=subbibliography,title={Internetquellen}]

% --- KI-Verzeichnis (PFLICHT seit 2024) ---
\newpage
% =============================================================================
% KI-Hilfsmittelverzeichnis
% -----------------------------------------------------------------------------
% PFLICHT, sobald KI-Werkzeuge bei der Erstellung der Arbeit genutzt wurden
% (FOM-Vorgaben zu KI in Studium und Lehre). Hier wird OFFENGELEGT, welche
% Werkzeuge wofür eingesetzt wurden. Diese Offenlegung ersetzt KEINE
% Quellenangabe: KI-gestützte Aussagen müssen zusätzlich mit Fachliteratur
% belegt werden.
%
% Trage jedes tatsächlich genutzte Werkzeug mit Anbieter, Version und
% Einsatzzweck ein. Die konkreten Prompts gehören (in Kurzform) in den Anhang.
% =============================================================================
\section*{KI-Hilfsmittelverzeichnis}
\addcontentsline{toc}{section}{KI-Hilfsmittelverzeichnis}

\begin{itemize}
  \item \textbf{[Werkzeug, Version (Anbieter)]:} [Einsatzzweck, z.\,B.
        Strukturierung der Gliederung, Recherche und Gegenprüfung von Quellen,
        sprachliche Überarbeitung einzelner Absätze].
  \item \textbf{[Werkzeug, Version (Anbieter)]:} [Einsatzzweck].
  \item \textbf{[Werkzeug, Version (Anbieter)]:} [Einsatzzweck].
\end{itemize}

Die durch KI erzeugten Ergebnisse wurden als Arbeitsgrundlage genutzt und
anschließend eigenständig geprüft, bewertet und überarbeitet. Die Verantwortung
für Inhalt, Quellenauswahl und Schlussfolgerungen liegt beim Verfasser.

% -----------------------------------------------------------------------------
% Hinweis: Prüfe in deinem aktuellen FOM-Leitfaden bzw. den Modulvorgaben die
% exakt geforderte Form (manche Standorte verlangen eine tabellarische
% Darstellung mit Spalten: Werkzeug | Zweck | betroffene Kapitel | Datum).
% -----------------------------------------------------------------------------


% --- Anhang ---
\newpage
% =============================================================================
% Anhang
% -----------------------------------------------------------------------------
% Hierhin gehören u.a.:
%   - KI-Prompts (Dokumentationspflicht bei KI-Nutzung) -- in Kurzform (WI)
%     oder als vollstaendiger Verlauf, wenn deine Pruefenden das verlangen.
%     Manche kuendigen bei fehlender Dokumentation "nicht bestanden" an.
%   - eigene Vorarbeiten (Expose, Vorstudie) benennen -- die Eigenstaendig-
%     keitserklaerung erfasst auch Uebernahmen aus eigenen Arbeiten
%   - ergänzende Tabellen und Abbildungen
%   - Interview-Transkripte, Fragebögen (bei empirischen Arbeiten)
% =============================================================================
\section*{Anhang}
\addcontentsline{toc}{section}{Anhang}

\subsection*{A. Dokumentation der KI-Prompts}

Im Folgenden werden die im Arbeitsprozess tatsächlich verwendeten Prompts in
Kurzform dokumentiert:

\begin{enumerate}
  \item \textbf{[Werkzeug]:} \enquote{[Kurzfassung des Prompts, z.\,B. Auftrag
        zur Strukturplanung der Gliederung]}
  \item \textbf{[Werkzeug]:} \enquote{[Kurzfassung des Prompts, z.\,B. Auftrag
        zur sprachlichen Überarbeitung eines Absatzes ohne inhaltliche
        Ergänzung]}
\end{enumerate}

% Längere Prompts kannst du als Codeblock einbinden:
% \begin{lstlisting}[breaklines=true,basicstyle=\ttfamily\scriptsize]
% Hier der wörtliche Prompt ...
% \end{lstlisting}


% --- Eigenständigkeitserklärung ---
% Manche Prüfungsämter verlangen die Erklärung digital im Upload-System,
% andere als Teil der PDF. Wenn sie ins Dokument soll: die zwei Zeilen
% unten entkommentieren (% entfernen).
% \newpage
% % =============================================================================
% Eigenständigkeitserklärung (FOM)
% =============================================================================
\section*{Eigenständigkeitserklärung}
\addcontentsline{toc}{section}{Eigenständigkeitserklärung}

Ich versichere, dass ich die vorliegende Arbeit selbständig und ohne unzulässige Hilfe
Dritter und ohne Benutzung anderer als der angegebenen Hilfsmittel angefertigt habe.

Die aus anderen Quellen direkt oder indirekt übernommenen Daten und Konzepte sind
unter Angabe der Quelle gekennzeichnet. Dies gilt auch für Quellen aus eigenen Arbeiten.

Ich versichere, dass ich alle KI-gestützten Hilfsmittel im KI-Hilfsmittelverzeichnis
vollständig aufgeführt habe und die mit ihrer Hilfe erstellten Inhalte im Text
entsprechend kenntlich gemacht habe.

\vspace{2cm}

\noindent
\begin{tabular}{@{}p{6cm}p{6cm}@{}}
  \hrulefill & \hrulefill \\
  Ort, Datum & Unterschrift \\
\end{tabular}


\end{document}
